% ===================== Conclusion and outlook (thesis Ch. 8) =====================
\section{Conclusions}

% Slide 49
\begin{frame}{Conclusions}
  \begin{columns}[t]
    \begin{column}{0.53\textwidth}
      \begin{block}{Results}
        \begin{itemize}
          \item The $x$-dependence of the Sivers shift can be computed from lattice QCD: scan $\bvec\tcdot P\neq0$ with off-axis staples, model with a Fourier-transformable power law, extrapolate $\zetahat\to\infty$.
          \item $\vec v_T\perp\vec b_T$ satisfies the kinematic constraint; symmetry improvement and interpolation give the off-axis staples.
          \item Gaussian ansatz excluded; a coupled power law describes all three amplitudes.
          \item At $|\vprp b|=0.34$ fm, $\langle k_y\rangle_{TU}$ peaks near $0.21$ GeV at small $x$; close to factorized in $x$ and $|\vprp b|$ (5.7--7.6\%).
        \end{itemize}
      \end{block}
    \end{column}
    \begin{column}{0.45\textwidth}
      \begin{alertblock}{Limitations}
        \begin{itemize}
          \item Bands statistical (Jackknife); model dependence explored, not propagated.
          \item Small $x$ depends on the fitted large-$\bvec\tcdot P$ tail.
          \item $\eta\to\infty$ through the asymptotic $c(\eta)$ model; $\zetahat\to\infty$ by extrapolation ($\zetahat\le0.79$).
          \item One ensemble: no continuum, volume or quark-mass study.
          \item $|\vprp b|\ge 3a$; $\tAmp_{12B}^{\text{Im}}$ set to zero.
        \end{itemize}
      \end{alertblock}
    \end{column}
  \end{columns}
\end{frame}

% Slide 50
\begin{frame}{Outlook}
  \quick
  \begin{itemize}
    \item The same method applies to the other leading-twist shifts. Each needs the amplitudes over a range of $\bvec\tcdot P$:
  \end{itemize}
  \vspace{1mm}
  \centering
  \renewcommand{\arraystretch}{1.25}
  \begin{tabular}{lll}
    \toprule
    Observable & Amplitude & Property \\
    \midrule
    Boer--Mulders shift $\langle k_y\rangle_{UT}(x)$ & $\tAmp_{4B}$ & T-odd: sign change SIDIS $\leftrightarrow$ DY \\
    Worm-gear shift $\langle k_x\rangle_{TL}(x)$ ($g_{1T}$) & $\tAmp_{7B}$ & T-even \\
    Generalized tensor charge $\tilde h_1^{(0)}/\tilde f_1^{(0)}(x)$ & $\tAmp_{9Bm}$ & T-even \\
    \bottomrule
  \end{tabular}
  \raggedright
  \begin{itemize}
    \item Larger momenta, more ensembles, a model-averaged systematic, and general $\vec v_T\cdot\vec b_T\neq0$ geometries.
  \end{itemize}
  \vspace{6mm}
  \centering
  {\Large \textbf{Thank you!}}
\end{frame}
